\fancyhead[L]{\small\textsf{\textit{CÁC NGUYÊN LÝ GIẢI BÀI TOÁN}}}
\fancyhead[R]{\small\textsf{\thepage}}

\noindent
\begin{minipage}[t]{0.26\textwidth}
    \vspace{8.04cm}
    \noindent
    \tcbox[
        colback=stewartpurple!15,
        colframe=stewartpurple!15,
        arc=0mm,
        top=1.5mm, bottom=1.5mm, left=2mm, right=2mm,
        boxrule=0pt
    ]{\textbf{\textsf{\color{stewartpurple}3\quad THỰC HIỆN KẾ HOẠCH}}}

    \vspace{0.61cm}
    \noindent
    \tcbox[
        colback=stewartpurple!15,
        colframe=stewartpurple!15,
        arc=0mm,
        top=1.5mm, bottom=1.5mm, left=2mm, right=2mm,
        boxrule=0pt
    ]{\textbf{\textsf{\color{stewartpurple}4\quad NHÌN LẠI}}}

    \vspace{3.77cm}
    \noindent
    \begin{tcolorbox}[
        colback=stewartpurple!15,
        colframe=stewartpurple!15,
        arc=0mm,
        top=1mm, bottom=1mm, left=1mm, right=2mm,
        boxrule=0pt
    ]
        \colorbox{stewartpurple}{\textbf{\textsf{\color{white}\footnotesize PS}}}\; \textbf{\textsf{\footnotesize Hiểu bài toán.}}
    \end{tcolorbox}
\end{minipage}%
\hfill
\begin{minipage}[t]{0.71\textwidth}
    \paragraph{\textsf{\textbf{\color{stewartcyan}Suy luận ngược}}} Đôi khi việc tưởng tượng rằng bài toán đã được giải quyết rồi đi ngược lại từng bước cho đến khi chạm tới các dữ kiện đã cho là rất hữu ích. Khi đó, bạn có thể đảo ngược các bước của mình và qua đó xây dựng lời giải cho bài toán ban đầu. Quy trình này thường được sử dụng trong việc giải phương trình. Chẳng hạn, khi giải phương trình $3x - 5 = 7$, ta giả sử $x$ là một số thỏa mãn $3x - 5 = 7$ và làm ngược lại. Ta cộng $5$ vào cả hai vế của phương trình rồi chia cả hai vế cho $3$ để được $x = 4$. Vì mỗi bước này đều có thể đảo ngược lại, nên ta đã giải xong bài toán.

    \paragraph{\textsf{\textbf{\color{stewartcyan}Thiết lập các mục tiêu phụ}}} Trong một bài toán phức tạp, việc đặt ra các mục tiêu phụ (trong đó tình huống mong muốn mới chỉ được đáp ứng một phần) thường rất hữu ích. Nếu trước tiên ta đạt được các mục tiêu phụ này, ta có thể dựa vào chúng để tiến tới mục tiêu cuối cùng.

    \paragraph{\textsf{\textbf{\color{stewartcyan}Lập luận gián tiếp}}} Đôi khi tiếp cận một bài toán một cách gián tiếp lại là phương pháp thích hợp. Khi sử dụng phép chứng minh phản chứng để chỉ ra $P$ suy ra $Q$, ta giả thiết rằng $P$ đúng và $Q$ sai, sau đó cố gắng tìm hiểu xem vì sao điều này không thể xảy ra. Bằng một cách nào đó, ta phải sử dụng thông tin này để dẫn tới một mâu thuẫn với điều mà ta chắc chắn là đúng.

    \paragraph{\textsf{\textbf{\color{stewartcyan}Quy nạp toán học}}} Khi chứng minh các mệnh đề có liên quan đến một số nguyên dương $n$, nguyên lý sau đây thường rất hữu ích.

    \vspace{0.16cm}
    \begin{tcolorbox}[
        colback=white,
        colframe=stewartred,
        arc=1mm,
        boxrule=1pt,
        top=3mm, bottom=3mm, left=4mm, right=4mm
    ]
        \textbf{\textsf{\color{stewartred}Nguyên lý quy nạp toán học}}\; Giả sử $S_n$ là một mệnh đề liên quan đến số nguyên dương $n$. Giả sử rằng:
        \begin{enumerate}
            \item[\textbf{1.}] $S_1$ đúng.
            \item[\textbf{2.}] $S_{k+1}$ đúng bất cứ khi nào $S_k$ đúng.
        \end{enumerate}
        Khi đó $S_n$ đúng với mọi số nguyên dương $n$.
    \end{tcolorbox}
    \vspace{0.16cm}

    Điều này là hợp lý vì do $S_1$ đúng, theo điều kiện 2 (với $k = 1$) suy ra $S_2$ đúng. Tiếp tục dùng điều kiện 2 với $k = 2$, ta thấy $S_3$ đúng. Lại dùng điều kiện 2 một lần nữa với $k = 3$, ta có $S_4$ đúng. Quy trình này có thể được tiếp tục vô hạn lần.

    \vspace{0.25cm}
    Ở Bước 2, một kế hoạch đã được vạch ra. Khi thực hiện kế hoạch đó, chúng ta phải kiểm tra từng giai đoạn và viết ra các chi tiết để chứng minh rằng mỗi bước đều chính xác.

    \vspace{0.25cm}
    Sau khi hoàn thành lời giải, ta nên xem lại toàn bộ bài làm, một phần để kiểm tra xem có sai sót hay không, và một phần để xem liệu ta có thể nghĩ ra cách giải nào dễ dàng hơn hay không. Một lý do khác của việc nhìn lại là nó giúp ta làm quen với phương pháp giải và điều này có thể hữu ích khi giải quyết các bài toán trong tương lai. Descartes từng nói: “Mỗi bài toán tôi giải được đều trở thành một quy tắc để sau đó dùng giải các bài toán khác.”

    \vspace{0.25cm}
    Các nguyên lý giải quyết vấn đề này được minh họa qua các ví dụ sau đây. Trước khi xem lời giải, bạn hãy thử tự mình giải các bài toán này, đồng thời tham khảo các nguyên lý giải quyết vấn đề nếu bị bế tắc. Bạn cũng có thể thấy hữu ích khi thỉnh thoảng quay lại tham khảo mục này trong quá trình làm bài tập ở các chương còn lại của cuốn sách.

    \vspace{0.33cm}
    \noindent
    \textbf{\textsf{\color{stewartcyan}VÍ DỤ 1}}\quad Biểu diễn cạnh huyền $h$ của một tam giác vuông có diện tích $25\text{ m}^2$ dưới dạng một hàm số theo chu vi $P$ của nó.

    \vspace{0.2cm}
    \noindent
    \textbf{\textsf{\color{stewartcyan}LỜI GIẢI}}\quad Trước hết, hãy sắp xếp các thông tin bằng cách xác định đại lượng chưa biết và các dữ kiện đã cho:
    \begin{align*}
        \textit{Chưa biết:} & \quad \text{cạnh huyền } h \\
        \textit{Các đại lượng đã cho:} & \quad \text{chu vi } P, \text{ diện tích } 25\text{ m}^2
    \end{align*}
\end{minipage}

\vfill
\begin{flushright}
    \small\textbf{71}
\end{flushright}